\documentclass[11pt]{article}
\usepackage[T1]{fontenc}
\usepackage{lmodern,amsmath,amssymb,amsthm,mathtools,booktabs,array,geometry,microtype}
\usepackage[colorlinks=true,linkcolor=blue,citecolor=blue,urlcolor=blue]{hyperref}
\geometry{margin=1in}
\newtheorem{theorem}{Theorem}[section]
\newtheorem{proposition}[theorem]{Proposition}
\newtheorem{lemma}[theorem]{Lemma}
\theoremstyle{remark}\newtheorem{remark}[theorem]{Remark}
\newcommand{\R}{\mathbb R}
\newcommand{\ev}{\mathrm{ev}}
\newcommand{\Per}{\operatorname{Per}}
\newcommand{\essinf}{\operatorname*{ess\,inf}}
\newcommand{\arc}{\operatorname{arc}}
\newcommand{\dd}{\,\mathrm d}
\newcommand{\norm}[1]{\lVert#1\rVert}
\title{P\'olya's inequality for the third planar Dirichlet eigenvalue}
\author{Liang Song}
\date{22 October 2026}
\hypersetup{pdftitle={Polya's inequality for the third planar Dirichlet eigenvalue},pdfauthor={Liang Song},pdfsubject={Spectral geometry; Dirichlet eigenvalues}}
\begin{document}
\maketitle
\begin{abstract}
We prove that $|\Omega|\lambda_3(\Omega)>12\pi$ for every planar open set of finite positive area, with the Dirichlet eigenvalues counted with multiplicity. No assumption on connectivity or boundary regularity is needed. The argument cuts the domain along an area bisector and reflects the two halves. Two estimates for cross-sectional profiles then give a dichotomy: either this cut has the required spectral bound, or a perpendicular cut does. The profile estimates use rearrangement and one-dimensional comparison, together with finite calculations in exact rational and interval arithmetic.
\end{abstract}
\section{Introduction}
Let $\lambda_k(\Omega)$ denote the Dirichlet eigenvalues of a planar open set, listed in increasing order and repeated according to multiplicity. P\'olya's conjecture asks whether
\begin{equation}\label{eq:polya}
 |\Omega|\lambda_k(\Omega)\ge 4\pi k
\end{equation}
for every $k\ge1$. P\'olya proved the conjecture for domains that tile the plane~\cite{Polya}. The first two eigenvalues satisfy stronger inequalities: the Faber--Krahn and Krahn--Szeg\H{o} bounds give $|\Omega|\lambda_1\ge\kappa$ and $|\Omega|\lambda_2\ge2\kappa$, where
\[
 \kappa=\pi j_{0,1}^2=18.1684\ldots
\]
and $j_{0,1}$ is the first positive zero of $J_0$; see~\cite{Henrot} for the classical bounds.

The third eigenvalue is the first index not settled by these two classical inequalities. They give only
\[
 |\Omega|\lambda_3(\Omega)\ge2\kappa=36.3368\ldots,
 \qquad 12\pi=37.6991\ldots.
\]
Our purpose is to cross the P\'olya threshold without restricting the topology of the domain.
\begin{theorem}\label{thm:main}
If $\Omega\subset\R^2$ is open and $0<|\Omega|<\infty$, then
\[
 |\Omega|\lambda_3(\Omega)>12\pi.
\]
\end{theorem}
Here and throughout, area means Lebesgue measure. The operator is defined by the quadratic form $u\mapsto\int_\Omega|\nabla u|^2$ on $H^1_0(\Omega)$. Thus holes contribute no area, all components are included, and no classical boundary trace is needed. The domain is allowed to be unbounded.

Results for every eigenvalue are known for disks and annuli~\cite{Balls,Annuli}. For general bounded planar open sets, a recent preprint of Speciel~\cite{Speciel} gives, among other estimates, $|\Omega|\lambda_k\ge(2/3)4\pi k$. Its strength is a bound valid at every index; at $k=3$ this particular corollary gives $8\pi$.
A preprint of Tang and Wang~\cite{TangWang} asserts the strict P\'olya inequality at every index for bounded simply connected planar domains. Theorem~\ref{thm:main} concerns one index but allows arbitrary topology, disconnected sets and unbounded sets of finite area. The proof below is independent of both preprints. Filling holes changes the area, while making cuts changes the Dirichlet form domain, so a simply connected result does not directly imply the present statement.

The constant in Theorem~\ref{thm:main} is the conjectured P\'olya constant, not a claimed optimal lower bound for $|\Omega|\lambda_3$. In particular, the theorem does not settle the separate problem of minimizing the third eigenvalue at fixed area.

The main geometric step is a choice of cutting line. Scale the area to one, bisect the set by a line, and reflect each half across it. Neumann bracketing reduces the problem to the second even eigenvalue on each reflected half. A small second even eigenvalue produces two disjoint low-energy pieces. Faber--Krahn forces each piece to have appreciable area, but this alone is insufficient: it gives no information about where that area lies. Our first profile estimate supplies the missing information. It forces each piece to occupy a definite amount of a strip next to the cutting line.

We rotate the bisecting line until the two adjacent strips have equal area. If this area is small, both reflected halves have large second even eigenvalue. If it is large, little area remains outside the full strip. A second profile estimate then works for a perpendicular bisector. These two alternatives cover every set.

Section~\ref{sec:geometry} carries out this geometric argument, assuming the two profile estimates. Sections~\ref{sec:one} and~\ref{sec:two} establish their analytic reductions. The remaining finite inequalities, and the procedures used to verify them, are given in Section~\ref{sec:cert}. The accompanying computational supplement contains the numerical certificate data and implementations of these procedures. No numerical approximation to an eigenvalue of a sample domain is used in the proof.
The dependencies can be read in the following order.
\begin{center}
\begin{tabular}{>{\raggedright\arraybackslash}p{.24\textwidth}>{\raggedright\arraybackslash}p{.64\textwidth}}\toprule
Step & Input and conclusion\\\midrule
Cut and reflect & Lemmas~\ref{lem:bracket}--\ref{lem:cores} reduce a small third eigenvalue to two disjoint low-energy pieces.\\
Single profile & Lemmas~\ref{lem:cap}, \ref{lem:strip-rearrangement}, \ref{lem:window-energy} and \ref{lem:bv-comparison}, with the neck, T and S certificates, prove Proposition~\ref{prop:single}.\\
Pair of profiles & Section~\ref{sec:two} and certificate B prove Proposition~\ref{prop:pair}.\\
Choice of direction & Lemma~\ref{lem:strip} and an area bisector with balanced adjacent strips yield Theorem~\ref{thm:main}.\\\bottomrule
\end{tabular}
\end{center}


\section{The geometric reduction}\label{sec:geometry}
We use the constants
\begin{equation}\label{eq:constants}
 \Lambda=12\pi,\quad \Lambda_0=\frac{377}{10},\quad
 \tau=\frac9{20},\quad m=\frac{1141}{5000},\quad a=\frac{519}{2000}.
\end{equation}
Thus $2\tau=0.9$, $1-4m=0.0872$, and $\Lambda<\Lambda_0$. All terminating decimals below denote exact rational numbers unless a certificate file is explicitly described as binary data.

\subsection{Cutting and reflecting}
For the moment let $|\Omega|=1$, and write the cutting line as $L=\{y=0\}$. Let $R_L$ be reflection across $L$, and set
\[
 \Omega_\pm=\Omega\cap\{\pm y>0\},\qquad
 S_\pm=\Omega_\pm\cup R_L\Omega_\pm\cup(\Omega\cap L).
\]
The sets $S_\pm$ are open. At a point of $\Omega\cap L$, an open ball in $\Omega$ supplies both half-balls after reflection. Let $\lambda_j^{\ev}(S_\pm)$ denote the eigenvalues restricted to functions even under $R_L$.

\begin{lemma}\label{lem:bracket}
For every line $L$,
\[
 \lambda_3(\Omega)\ge
 \min\{\lambda_2^{\ev}(S_+),\lambda_2^{\ev}(S_-)\}.
\]
\end{lemma}
\begin{proof}
Let $V_\pm$ be the closure in $H^1(\Omega_\pm)$ of restrictions of $C_c^\infty(\Omega)$ functions. Restriction to the two halves preserves the sum of the $L^2$ norms and the sum of the Dirichlet energies. The min--max principle bounds the eigenvalues of $\Omega$ from below by those of $V_+\oplus V_-$. Among its first three eigenvalues, at least two come from the same summand. Its third eigenvalue is therefore at least the smaller of the second eigenvalues of the two summands.

Even reflection maps $V_\pm$ into $H^1_0(S_\pm)^{\ev}$ and multiplies both energy and squared norm by two. This follows first for the restricted smooth functions and then by closure. Another application of min--max proves the claim.
\end{proof}

\begin{lemma}[Finite-area spectral setting]\label{lem:finite-area}
For every open set $U\subset\R^2$ of finite positive area, the Dirichlet form on $H^1_0(U)$ has a discrete positive spectrum, with finite multiplicities and the usual min--max characterization.
\end{lemma}
\begin{proof}
Extend functions by zero to $\R^2$. Faber--Krahn in its support-measure form gives
\[
 \norm u_2^2\le\frac{|U|}{\kappa}\norm{\nabla u}_2^2
 \qquad(u\in H^1_0(U)),
\]
first on smooth compactly supported functions and then by density. This makes the Dirichlet form coercive. To prove compactness, let $(u_n)$ be bounded in $H^1_0(U)$ and fix $p>2$. The planar Sobolev inequality and H\"older give
\[
 \int_{U\setminus B_R}|u_n|^2
 \le |U\setminus B_R|^{1-2/p}\norm{u_n}_{L^p(\R^2)}^2
 \le C|U\setminus B_R|^{1-2/p}.
\]
The right-hand side tends to zero uniformly in $n$. Rellich compactness on balls, followed by a diagonal subsequence and this tail bound, gives convergence in $L^2(U)$. The closed-form construction and the spectral theorem for its compact inverse now give the assertion. All multiplicities are retained. In particular, boundedness or boundary regularity of $U$ is not used.
\end{proof}


\begin{lemma}\label{lem:cores}
There are disjoint open sets $W_1,W_2\subset S_+$, both invariant under $R_L$, such that
\[
 \lambda_1(W_i)\le\lambda_2^{\ev}(S_+),\qquad i=1,2.
\]
\end{lemma}
\begin{proof}
Even eigenfunctions are genuine weak eigenfunctions on $S_+$: the weak equation holds for even test functions, and both sides vanish for odd ones. They are smooth in the interior.

Choose a nonnegative first even eigenfunction $u_1$ and a second even eigenfunction $u_2$ orthogonal to it. If $u_2$ changes sign, use $W_1=\{u_2>0\}$ and $W_2=\{u_2<0\}$. If it has one sign, change its sign if necessary. Orthogonality and continuity imply that $\{u_1>0\}$ and $\{u_2>0\}$ are disjoint, and we use these sets. This also applies when the first even eigenvalue is multiple: absolute value preserves the Rayleigh quotient of a minimizer, so a nonnegative minimizer can still be chosen.

The sign parts are admissible Dirichlet test functions on their positivity sets. To see this, if $u\in H^1_0(U)\cap C(U)$, approximate $u$ in $H^1$ by $u_n\in C_c^\infty(U)$. The functions $\min\{u_n^+,(u-\varepsilon)^+\}$ have compact support in $\{u>0\}$, are bounded in $H^1$, and converge in $L^2$ to $(u-\varepsilon)^+$. Weak closure, followed by $\varepsilon\downarrow0$, gives $u^+\in H^1_0(\{u>0\})$. Testing the eigenvalue equation with the sign parts gives the Rayleigh bounds. Evenness gives the required symmetry.
\end{proof}

\subsection{The two profile estimates}
For an even nonnegative profile $c$, define
\[
 S_c=\{(x,y):|x|<c(y)/2\}.
\]
When $c$ is lower semicontinuous this is open. For a measurable profile define
\[
 \mu[c]=\inf_{F\ne0}
 \frac{\int_\R\bigl(|F'|^2+\pi^2|F|^2/c^2\bigr)}{\int_\R|F|^2}.
\]
The infimum is over $H^1(\R)$ functions with finite numerator; the potential is infinite on $\{c=0\}$ and zero where $c=\infty$.

We will prove the following two statements.
\begin{proposition}[A single profile]\label{prop:single}
If $c$ is even, lower semicontinuous and nonnegative, and
\[
 \int_0^\infty c\le a,\qquad \int_0^\tau c\le m,
\]
then $\lambda_1(S_c)>\Lambda$.
\end{proposition}
\begin{proposition}[A pair of profiles]\label{prop:pair}
Suppose $c_1,c_2$ are even nonnegative measurable profiles and
\[
 \int_0^\infty(c_1+c_2)\le\frac12,\qquad
 \int_0^\infty(c_1+c_2-0.9)_+\le0.0872.
\]
Then $\max\{\mu[c_1],\mu[c_2]\}\ge\Lambda_0$.
\end{proposition}
The first estimate rules out a low-energy piece with too little area near the line. The second rules out two low-energy pieces whose combined width exceeds $0.9$ only by a small total amount. The constants have been chosen so that the two estimates fit the geometric alternatives: $4m$ is the area forced into a full strip, and $1-4m$ is the area left outside it.

Write $W^y=\{x:(x,y)\in W\}$ and $c_W(y)=|W^y|$. For an open set $W$, Fatou's lemma shows that $c_W$ is lower semicontinuous. Steiner rearrangement gives
\begin{equation}\label{eq:steiner}
 \lambda_1(S_{c_W})\le\lambda_1(W).
\end{equation}
There is also a one-dimensional lower bound:
\begin{equation}\label{eq:slicing}
 \mu[c_W]\le\lambda_1(W).
\end{equation}
On each component of $W^y$, the Dirichlet inequality gives $\int|u_x|^2\ge\pi^2c_W(y)^{-2}\int|u|^2$. For $F(y)=\norm{u(\cdot,y)}_2$, one has $|F'|\le\norm{u_y(\cdot,y)}_2$ almost everywhere. Integration proves~\eqref{eq:slicing}. Both arguments extend from smooth compactly supported functions by approximation. Rearrangement is used in~\eqref{eq:steiner}; slicing is used in~\eqref{eq:slicing}. They serve different purposes.

\begin{lemma}\label{lem:strip}
Suppose $L$ bisects the area of $\Omega$. If $\lambda_2^{\ev}(S_+)\le\Lambda$, then
\[
 |\Omega\cap\{0<y<\tau\}|>2m.
\]
\end{lemma}
\begin{proof}
Take the sets from Lemma~\ref{lem:cores} and put $c_i=c_{W_i}$ and $a_i=\int_0^\infty c_i$. Faber--Krahn and disjointness give
\[
 \frac{\kappa}{2\Lambda}\le a_i,
 \qquad a_1+a_2\le\frac12.
\]
Consequently $a_i\le\frac12-\kappa/(2\Lambda)<a$. By~\eqref{eq:steiner}, $\lambda_1(S_{c_i})\le\Lambda$. Proposition~\ref{prop:single} forces $\int_0^\tau c_i>m$ for each $i$. Add these inequalities and use $c_1+c_2\le c_{\Omega_+}$.
\end{proof}

\subsection{Choosing the line}
\begin{proof}[Proof of Theorem~\ref{thm:main}, assuming Propositions~\ref{prop:single} and~\ref{prop:pair}]
Continue to assume $|\Omega|=1$. For a unit vector $n$, an area median is a number $s$ with $|\Omega\cap\{x\cdot n<s\}|=1/2$.

First suppose a direction has an interval of medians. The open strip between its extreme medians is disjoint from $\Omega$: it has zero intersection area, and a nonempty open intersection would have positive area. Cut strictly inside this strip. Each reflected half is a disjoint union of two copies of the original half. If its second even eigenvalue were at most $\Lambda$, Lemma~\ref{lem:cores} and Faber--Krahn, applied to the upper halves of the two pieces, would give
\[
 \frac12\ge |W_1^+|+|W_2^+|\ge\frac{2\kappa}{\Lambda}>\frac12,
\]
a contradiction. Lemma~\ref{lem:bracket} finishes this case.

We may assume that every median $s(n)$ is unique. The function $(n,s)\mapsto|\Omega\cap\{x\cdot n<s\}|$ is continuous by dominated convergence, since a line has planar measure zero. Its unique median depends continuously on $n$: choose two levels on either side of a given median at which the areas lie strictly below and above $1/2$; continuity traps nearby medians between them.

Set
\[
 q(n)=|\Omega\cap\{s(n)<x\cdot n<s(n)+\tau\}|.
\]
This is continuous, and $s(-n)=-s(n)$. The continuous odd function $q(n)-q(-n)$ vanishes somewhere on the circle. Rotate coordinates so this median line is $y=0$, and write $Q=q(n)=q(-n)$.

If $Q\le2m$, Lemma~\ref{lem:strip} shows that both second even eigenvalues exceed $\Lambda$. If $Q>2m$, the full horizontal strip has area $2Q>4m$. Choose a vertical area bisector. At horizontal position $x$, let $e(x)$ be the length of the part of the vertical section outside $|y|<\tau$. Its full length is at most $0.9+e(x)$, and
\[
 \int_\R e(x)\dd x<1-4m=0.0872.
\]
On either side of the vertical bisector, the full profile has mass $1/2$ and excess above $0.9$ less than $0.0872$. If that reflected half had second even eigenvalue below $\Lambda_0$, its two pieces from Lemma~\ref{lem:cores} would, by~\eqref{eq:slicing}, contradict Proposition~\ref{prop:pair}. Both second even eigenvalues are therefore at least $\Lambda_0>\Lambda$.

Lemma~\ref{lem:bracket} gives $\lambda_3(\Omega)>\Lambda$ in both alternatives. Finally $\lambda_k(t\Omega)=t^{-2}\lambda_k(\Omega)$ removes the area normalization.
\end{proof}

\section{The single-profile estimate}\label{sec:one}
We prove Proposition~\ref{prop:single}. Suppose, to obtain a contradiction, that
\[
 \lambda=\lambda_1(S_c)\le\Lambda,
 \qquad m'=\int_0^\tau c\le m.
\]
Let $T=\R\times(-\tau,\tau)$, $D=S_c\cap T$, and $\alpha=a-m'$. There are two sources of energy to exploit. Very little area in the central strip creates a narrow neck. Otherwise, the two parts outside the strip force a boundary penalty on its walls. After replacing these outer parts by that penalty, only the size of the nonzero boundary trace remains to be considered.

\subsection{A narrow central strip}
Only the component of $S_c$ that meets $y=0$ needs consideration. Every other component lies in a half-plane and has area at most $a$, so its first eigenvalue is at least $\kappa/a>\Lambda$. The component meeting the axis is unique and symmetric; its positive ground state is even.

Let $U_t$ be a superlevel set of that ground state, of area $s$. Some good horizontal slice $y_t\in(0,\tau)$ has length at most $m_0/\tau$, where $m_0=0.03$. This follows by averaging the section lengths. The part of $U_t$ above that slice has area at least $s/2-m_0$. The isoperimetric inequality, applied to this part and then reflected, gives
\begin{equation}\label{eq:neck}
 \Per(U_t)\ge I(s):=
 \max\left\{\sqrt{4\pi s},\,
 2\left[\sqrt{4\pi(s/2-m_0)_+}-\frac{m_0}{\tau}\right]\right\}.
\end{equation}
The subtracted term accounts for the new cut. The same argument for finite-perimeter superlevel sets follows from the slicing and gluing formulae for perimeter.

The coarea formula and Cauchy--Schwarz give
\begin{equation}\label{eq:neck-quotient}
 \lambda\ge\inf_{g(2a)=0}
 \frac{\int_0^{2a}I(s)^2|g'(s)|^2\dd s}{\int_0^{2a}|g(s)|^2\dd s}.
\end{equation}
Here $g$ is the decreasing rearrangement in the area variable, extended by zero if the support has area less than $2a$. This is the energy form of the rearrangement argument~\cite{Talenti,BrothersZiemer}.

The certificate in Section~\ref{sec:cert} bounds the quotient in~\eqref{eq:neck-quotient} from below by $45.8318$. Its comparison principle is short. If $\zeta(0)=0$, $\zeta$ remains finite, and
\[
 \zeta'\ge L+\frac{\zeta^2}{I^2},
\]
with only upward jumps, then completion of the square and integration by parts yield
\begin{align*}
 \int I^2g'^2-L\int g^2
 &=\int I^2\left(g'+\frac{\zeta g}{I^2}\right)^2\\
 &\quad+\int g^2\left(\dd\zeta-\frac{\zeta^2}{I^2}\dd s-L\dd s\right)
   -[\zeta g^2]\ge0.
\end{align*}
The formula is first applied away from zero and then passed to the limit. The endpoint terms vanish. Taking $L=45.8318$ contradicts $\lambda\le\Lambda$.

\subsection{Removing the caps}
Assume $m'\ge0.03$. The following minimization supplies the Robin coefficient.
\begin{lemma}\label{lem:cap}
For a nonnegative width $b$ with $\int_0^\infty b\le\alpha$, and an admissible function $G$ with $G(0)=1$,
\[
 \int_0^\infty\left(G'^2+\frac{\pi^2G^2}{b^2}\right)\dd y
 \ge\frac{\sigma_1}{\sqrt\alpha},\qquad
 \sigma_1=\frac32\left(\frac\pi{\sqrt3}\right)^{1/2}.
\]
The constant is sharp.
\end{lemma}
\begin{proof}
Scaling reduces the claim to $\alpha=1$. We may replace $G$ by $|G|$. Put $B=\int G^{2/3}$. H\"older's inequality gives $\int G^2/b^2\ge B^3$. Also, integration of $(G^{4/3})'$ followed by Cauchy--Schwarz gives
\[
 1\le\frac43\left(\int G'^2\right)^{1/2}B^{1/2}.
\]
Thus the energy is at least $9/(16B)+\pi^2B^3$. Its minimum occurs at $B^4=3/(16\pi^2)$ and equals $\sigma_1$. Equality is attained by $G(y)=(1-y/Y)_+^{3/2}$ with $Y^4=3/\pi^2$, and by the corresponding width proportional to $G^{2/3}$.
\end{proof}

Write $u$ for a ground state on $S_c$. Its energy on either cap is at least $\sigma_1\alpha^{-1/2}$ times the squared $L^2$ trace on the adjoining wall, by slicing and Lemma~\ref{lem:cap}. Even reflection of a cap function across the wall gives a whole-plane $H^1$ function supported on a set of area at most $2\alpha$. Faber--Krahn in its support-measure form therefore gives
\[
 M_C\le\frac{E_C}{\nu},\qquad \nu=\frac\kappa{2\alpha},
\]
for the total cap mass and energy as well. Throughout $0.03\le m'\le m$, one has $\nu>\Lambda_0$.

Let $M$ and $I$ be the mass and energy inside $T$. For either $\widetilde\Lambda=\Lambda$ or $\Lambda_0$, the identity $\lambda(M+M_C)=I+E_C$ implies
\[
 \lambda M\ge I+\left(1-\frac{\widetilde\Lambda}{\nu}\right)E_C
 \ge I+\sigma_{\widetilde\Lambda}(m')\int_{\partial T}u^2,
\]
where
\begin{equation}\label{eq:sigma}
 \sigma_{\widetilde\Lambda}(m')=
 \frac{\sigma_1}{\sqrt{a-m'}}
 \left(1-\frac{2\widetilde\Lambda(a-m')}{\kappa}\right)>0.
\end{equation}
We used $\lambda\le\Lambda\le\widetilde\Lambda$. If $M=0$, the cap estimate gives $\lambda\ge\nu$, already a contradiction.

For arbitrary boundaries it is convenient to work with
\[
 \mathcal V(D)=\{v\in H^1(T):v=0\text{ a.e. on }T\setminus D\},\quad
 Q_\sigma(v)=\int_T|\nabla v|^2+\sigma\int_{\partial T}v^2.
\]
Zero extension of $u$ belongs to $H^1(\R^2)$, so $u|_T\in\mathcal V(D)$. We have proved
\begin{equation}\label{eq:robin}
 \inf_{v\in\mathcal V(D)\setminus\{0\}}
 \frac{Q_{\sigma_{\widetilde\Lambda}(m')}(v)}{\norm v_2^2}\le\lambda.
\end{equation}

\subsection{Rearrangement and the boundary window}
Take $v\in\mathcal V(D)$. First replace it by $|v|$, then by
$\bigl(v(x,y)^2+v(x,-y)^2\bigr)^{1/2}/\sqrt2$.
These operations preserve mass and the boundary term and do not increase energy. Rearrange each positive vertical half-fibre decreasingly toward $y=0$, and reflect evenly. The trace at $y=\tau$ becomes the essential infimum on that half-fibre, so its square integral cannot increase. Finally perform horizontal Steiner rearrangement.

These operations preserve the $L^2$ norm and decrease Dirichlet energy. The one-sided rearrangement is performed on a finite interval; it imposes no zero boundary value at its outer endpoint. Horizontal rearrangement can be obtained by polarizations~\cite{BrockSolynin}. These exchange points at the same height and thus apply on a strip. The trace commutes with this rearrangement because $y\mapsto v(\cdot,y)$ is continuous into $L^2(\R)$ and rearrangement is an $L^2$ contraction.

The resulting function is nonnegative and decreases with both $|x|$ and $|y|$. It is supported in the doubly symmetric set
\[
 D^*=\{(x,y):|y|<|D^x|/2\},\qquad |D^*|=2m'.
\]
Put $w_* = \essinf_{|y|<\tau}c(y)$. The columns of $D^*$ that fill the entire strip form, up to endpoints, an interval of width $w_*$. Hence $2\tau w_*\le2m'$, and the wall traces vanish outside that interval. If $w_*=0$, extension by zero across the walls and Faber--Krahn give
\[
 Q_\sigma(v)\ge\frac\kappa{2m'}\norm v_2^2
 \ge\frac\kappa{2m}\norm v_2^2>\Lambda_0\norm v_2^2.
\]
It remains to consider a positive boundary window.

\subsection{Small windows: a perimeter comparison}
For a chord of length $2\tau$, let $\arc(2\tau,A)$ be the length of the circular arc bounding a segment of area $A$ on one side of the chord. If a doubly symmetric set $U\subset T$ has area $s$ and full-column width $\ell>0$, then
\begin{equation}\label{eq:perimeter}
 \Per(U;T)\ge P(\ell,s):=
 \max\left\{2\arc\left(2\tau,\frac{s-2\tau\ell}{2}\right),
                  \sqrt{4\pi s}-2\ell\right\}.
\end{equation}
The second bound follows by adding the two wall traces to the relative perimeter and applying the isoperimetric inequality. For the first, consider the part to the right of the full columns. Complete it on the other side of the chord by the complementary circular segment. Choose the circle so that the completed set has the same area as the circle. The isoperimetric inequality says that its unknown boundary is no shorter than the missing circular arc. Reflecting gives the factor two. The perimeter includes any vertical boundary left on the chord, so no endpoint segment is discarded.

Apply~\eqref{eq:perimeter} to the level sets of a doubly decreasing $v$. Write $q(s)$ for its decreasing rearrangement in the area variable and $\ell(s)$ for the corresponding wall-trace width. Then $\ell$ is nondecreasing, $0\le\ell\le w_*$ and $2\tau\ell(s)\le s$. Coarea and the layer-cake formula give
\begin{equation}\label{eq:window-energy}
 \begin{split}
 Q_\sigma(v)&\ge\int_0^{A_0}P(\ell(s),s)^2q'(s)^2\dd s
                   +2\sigma\int_{[0,A_0]}q(s)^2\dd\ell(s),\\
 \norm v_2^2&=\int_0^{A_0}q(s)^2\dd s,
 \end{split}
\end{equation}
where $A_0\ge|D^*|$ and $q$ is extended by zero. The Stieltjes measure includes jumps; its convention is $\ell(0^-)=0$. At a plateau of $q$, the jump is placed at the right endpoint. This is the convention obtained from the wall layer-cake formula. The full-column width and the positive trace width agree for almost every level; exceptional constant levels do not alter coarea.

The relative perimeter is bounded below by $\min\{\sqrt{4\pi s},4\tau\}>0$ for $s>0$. Thus coarea gives local absolute continuity of $q$ away from zero. One first truncates $v$ at a finite height, so $q$ is bounded, and removes the truncation at the end.

Here is the criterion used for~\eqref{eq:window-energy}. Let $\widetilde P$ be a lower bound for $P$ on a rectangle $0\le s\le A_0$, $0\le\ell\le w$. On $\ell=0$ use $\widetilde P=\sqrt{4\pi s}$. Suppose a continuous, piecewise smooth function $\rho$ satisfies
\begin{equation}\label{eq:rho}
 \rho(0,0)=0,\qquad
 \rho_s\le-\Lambda_0-\frac{\rho^2}{\widetilde P^2},\qquad
 \rho_\ell\le2\sigma.
\end{equation}
Along a nondecreasing path put $g(s)=\rho(s,\ell(s))$. Fill each jump by its vertical segment. The chain rule on the smooth pieces, and integration across these segments, give
\[
 \dd g\le-\left(\Lambda_0+\frac{g^2}{\widetilde P^2}\right)\dd s
            +2\sigma\dd\ell.
\]
Use $\widetilde P^2q'^2\ge2gqq'-g^2q^2/\widetilde P^2$ and integrate by parts. The result is
\[
 \int\widetilde P^2q'^2+2\sigma\int q^2\dd\ell
       -\Lambda_0\int q^2\ge0.
\]
The endpoint at $A_0$ vanishes since $q(A_0)=0$. This also holds when $A_0$ equals the support area: the perimeter lower bound is positive near that endpoint, so the energy estimate prevents a nonzero jump to the zero extension. For the paths in~\eqref{eq:window-energy}, $\ell(s)\le s/(2\tau)$; local Lipschitz continuity of $\rho$ and $\rho(0,0)=0$ therefore give $\rho(s,\ell(s))=O(s)$. Truncating at $s=\varepsilon$ and passing to the limit handles the origin.

No absolute continuity of $\ell$ is assumed here. On a compact subinterval away from zero, complete each jump of its graph by a vertical segment and approximate this monotone graph by polygonal graphs. The differential inequality holds along every segment. Uniform continuity of $\rho$ and continuity of the truncated $q$ pass the boundary integrals to their Stieltjes limits. This includes the singular continuous part of $\dd\ell$, as well as its atoms. On triangulation edges the tangential derivatives agree, since the adjacent pieces of $\rho$ are continuous there. Thus the comparison applies to the paths supplied by coarea, not only to smooth paths.

The finite certificates verify~\eqref{eq:rho} for all $0.03\le m'\le m$ and $w_*\le0.241$, with $\sigma=\sigma_{\Lambda_0}(m')$. Together with~\eqref{eq:robin}, this contradicts $\lambda\le\Lambda$.

\subsection{Large windows: a one-dimensional comparison}
For $w_*\ge0.240$, horizontal slicing is more efficient. The horizontal profile $c^*$ of $D^*$ satisfies
\[
 c^*(y)\ge0.240,\qquad \int_0^\tau c^*(y)\dd y=m'.
\]
The norm-of-slices argument in~\eqref{eq:slicing} gives
\[
 Q_\sigma(v)\ge \mu_h[c^*]\norm v_2^2,
\quad
 \mu_h[b]=\inf_{f\ne0}
 \frac{\int_0^\tau(f'^2+\pi^2f^2/b^2)\dd y+\sigma f(\tau)^2}
      {\int_0^\tau f^2\dd y}.
\]
The endpoint at zero is Neumann. No monotonicity assumption on $b$ is needed.

Let $F''=(\pi^2/b^2-\Lambda_0)F$, $F(0)=1$, $F'(0)=0$, and write $(F,F')=r(\sin\theta,\cos\theta)$ with $\theta(0)=\pi/2$. Then
\begin{equation}\label{eq:prufer-S}
 \theta'=f(\theta,b):=1+\sin^2\theta\left(\Lambda_0-1-\frac{\pi^2}{b^2}\right).
\end{equation}
The condition $\mu_h[b]\le\Lambda_0$ is equivalent to $\theta(\tau)\ge\theta_R:=\pi/2+\arctan\sigma$. Before a zero of $F$, this follows from
\[
 Q(\varphi)-\Lambda_0\int\varphi^2
 =\int F^2\bigl((\varphi/F)'\bigr)^2
   +\left(\sigma+\frac{F'(\tau)}{F(\tau)}\right)\varphi(\tau)^2.
\]
If $F$ has a zero, Sturm comparison gives the same conclusion. The angle cannot cross a multiple of $\pi$ downwards, since its derivative there is one.

A potential function bounds the mass of any trajectory satisfying this terminal condition. If a piecewise constant slope $p(\theta)$ satisfies
\begin{equation}\label{eq:potential}
 b-\nu-p(\theta)f(\theta,b)\ge0\qquad(b\ge w),
\end{equation}
then on an allowed trajectory segment,
\[
 \int b\dd y\ge \nu\,\Delta y+\int_{\theta_{\rm start}}^{\theta_{\rm end}}p(t)\dd t.
\]
This remains valid on time intervals spent at a slope discontinuity: there $\theta'=0$ almost everywhere and~\eqref{eq:potential} gives $b\ge\nu$.

To check~\eqref{eq:potential} on an angular cell, put $s=\sin^2\theta$ and minimize
\[
 G(p,s)=\min_{b\ge w}\left(b+\frac{p\pi^2s}{b^2}\right)
                 -\nu-p\bigl(1+(\Lambda_0-1)s\bigr).
\]
This is concave in $s$, being an infimum of affine functions. It suffices to check the extreme values of $s$ on the cell. A cell meeting $\pi/2$ must include $s=1$ in this test. The scalar minimum is explicit: with $K=p\pi^2s$, it is $w+K/w^2$ when $K\le0$ or $(2K)^{1/3}\le w$, and otherwise $\frac32(2K)^{1/3}$.

Appendix~\ref{app:stairs} explains how these potentials cover every trajectory, including those that turn around. The eight interval certificates give a mass lower bound strictly greater than $m'$ whenever $\mu_h[c^*]\le\Lambda_0$, with $\sigma=\sigma_\Lambda(m')$. This proves the large-window case. The two window ranges overlap. The range $m'<0.108$ is empty in this case because $0.240\tau=0.108$.

We have covered $m'\le0.03$, $w_*=0$, small positive windows, and large windows. Proposition~\ref{prop:single} follows. Its conclusion is only $\lambda_1(S_c)>12\pi$: the large-window Robin coefficient was derived under $\lambda\le12\pi$. The internal certificate value $37.7$ does not strengthen that conclusion to $\lambda_1(S_c)\ge37.7$.

\section{The pair-profile estimate}\label{sec:two}
We prove Proposition~\ref{prop:pair} by showing that two profiles with $\mu[c_i]<\Lambda_0$ and total half-line mass at most $1/2$ must have excess above $0.9$ greater than $0.0872$.

For positive bounded profiles on a finite interval, use $t=\sqrt{\Lambda_0}\,y$, $\ell_i=\sqrt{\Lambda_0}\,c_i$ and the scaled Pr\"ufer angles
\[
 \theta_i'=1-\frac{\pi^2\sin^2\theta_i}{\ell_i^2},\qquad \theta_i(0)=\pi/2.
\]
Sturm theory~\cite{Zettl} says that $\mu[c_i]<\Lambda_0$ forces each angle to reach $\pi$ in finite time. General profiles follow by truncating to an interval containing compactly supported test functions, capping large widths, and adding a small positive width. The strict Rayleigh inequalities survive these changes; the extra mass and excess tend to zero. The certificate below has strict slack in both the initial state and final excess, so this approximation does not lose the conclusion.

Set $C=3\sqrt3\pi/2$, $U_i=1+\cos\theta_i$ and $X_i=\sin\theta_i$. Before the first angle reaches $\pi$, $0<\theta_i<\pi$, so $X_i\ge0$. The elementary inequality
\begin{equation}\label{eq:angle-budget}
 \ell_i\ge C X_i\theta_i'
\end{equation}
follows, on writing $\ell_i=\sqrt3\pi z$, from $2z+X_i^3/z^2\ge3X_i$; the difference is $(z-X_i)^2(2z+X_i)/z^2$.

Let $M_*$ be the allowed physical half-line mass, initially $1/2$, and let $M(t)$ be the mass already used. Introduce
\[
 s(t)=\frac{\Lambda_0}{C}\bigl(M_*-M(t)\bigr)-U_1(t)-U_2(t).
\]
By~\eqref{eq:angle-budget}, $s$ is nonincreasing. It is nonnegative along a feasible trajectory: applying~\eqref{eq:angle-budget} to each angle's remaining journey to $\pi$ bounds the remaining mass below by $C(U_1+U_2)/\Lambda_0$. At the start,
\[
 s(0)=\frac{\Lambda_0}{3\sqrt3\pi}-2.
\]
Thus the trajectory lies in $[0,2]^2\times[0,s_{\rm top}]$, where the certificate chooses $s_{\rm top}$ slightly larger than this initial value.

Let $w(U_1,U_2,s)$ be continuous and affine on each tetrahedron of a triangulation. Write its gradient on a tetrahedron as $(q_1,q_2,-k)$, and put $P_i=q_i+k$. Require $k,P_1,P_2\ge0$ and $w\le0$ on each face $U_i=0$. Put
\[
 \omega=\frac{\sqrt{\Lambda_0}\,0.9}{\sqrt3\pi},\qquad
 A=P_1^{1/3}X_1+P_2^{1/3}X_2.
\]
The remaining cell condition is
\begin{equation}\label{eq:HJB}
 3(P_1X_1+P_2X_2)\le\Psi_k(A),
\end{equation}
where
\[
\Psi_k(A)=
\begin{cases}
3k^{2/3}A,&A\le k^{1/3}\omega,\\
2k\omega+A^3/\omega^2,&k^{1/3}\omega\le A\le(k+1/2)^{1/3}\omega,\\
3(k+1/2)^{2/3}A-\omega,&A\ge(k+1/2)^{1/3}\omega.
\end{cases}
\]
To see why this is sufficient, write $\ell_i=\sqrt3\pi z_i$ and $Z=z_1+z_2$. Direct differentiation shows that the required inequality along a trajectory is
\[
 \sum_i P_i\left(3X_i-\frac{X_i^3}{z_i^2}\right)
 \le 2kZ+(Z-\omega)_+.
\]
At fixed $Z$, H\"older's inequality gives
\[
 \inf_{z_1+z_2=Z}\sum_i\frac{P_iX_i^3}{z_i^2}=\frac{A^3}{Z^2}.
\]
Minimizing $A^3/Z^2+2kZ+(Z-\omega)_+$ over $Z>0$ gives exactly the three branches of $\Psi_k$. Thus~\eqref{eq:HJB} proves, for every control,
\[
 \frac{\dd}{\dd t}\left(\frac{2C}{\Lambda_0}w\right)
 \ge-\frac{(\ell_1+\ell_2-\sqrt{\Lambda_0}\,0.9)_+}{\Lambda_0}.
\]
The chain rule also holds when the trajectory follows a cell face: continuity of the affine pieces makes their tangential derivatives agree.

Stop when the first angle reaches $\pi$. The terminal value of $w$ is nonpositive. Since $w_s=-k\le0$, its initial value is at least $w(1,1,s_{\rm top})$. Integration gives
\[
 \int_0^\infty(c_1+c_2-0.9)_+\dd y
 \ge\frac{3\sqrt3\pi}{\Lambda_0}w(1,1,s_{\rm top}).
\]
The rational checker verifies that this lower bound exceeds $0.0879293$. This is strictly larger than $0.0872$, proving Proposition~\ref{prop:pair}.

\section{The finite certificates}\label{sec:cert}
It remains to verify a finite collection of inequalities. We describe the objects being certified, the exact tests applied to them, and the resulting bounds. We label the four calculations B (pair profiles), T (small windows), S (large windows), and N (the neck). The certificate arrays are supplied as supplementary data. Their construction plays no part in verification: the tests use the fixed arrays and do not run an optimizer.

\begin{center}
\begin{tabular}{>{\raggedright\arraybackslash}p{.06\textwidth}>{\raggedright\arraybackslash}p{.31\textwidth}>{\raggedright\arraybackslash}p{.49\textwidth}}\toprule
Test & Range & Certified condition and use\\\midrule
B & Combined half-area $\le1/2$, width threshold $.9$ & Excess mass $>.0879293>.0872$; Proposition~\ref{prop:pair}.\\
T & $m'\in[.03,.2282]$, $w_*\le.241$; nine boxes & Both inequalities in~\eqref{eq:rho} on all 192000 triangles per box; Lemma~\ref{lem:bv-comparison}.\\
S & $m'\in[.108,.2282]$, $w_*\ge.240$; eight intervals & Required mass exceeds the upper endpoint of each interval; the table below records the bounds.\\
N & $m'\le.03$, support area $\le.519$ & Riccati comparison at $L=45.8318>12\pi$; inequality~\eqref{eq:neck-quotient}.\\\bottomrule
\end{tabular}
\end{center}

\subsection{The pair-profile certificate}
The pair-profile certificate is a continuous piecewise affine function on 8664 Kuhn tetrahedra. The underlying box grid has $38\times38\times1$ cells, with $U_1,U_2\in[0,2]$ and
\[
 s_{\rm top}=\frac{2796334405517737}{9007199254740992}.
\]
Each box is divided into its six coordinate-order tetrahedra. The nodal values are first adjusted by subtracting $10^{-4}(U_1+U_2+s)$ and rounded down to multiples of $2^{-44}$. Values on a terminal face are replaced by their minimum with zero. These adjusted values define $w$. The finite test verifies the signs of $k,P_i$, both terminal faces, and the initial-state bound; it also checks $s_{\rm top}>\Lambda_0/(3\sqrt3\pi)-2$ with rational bounds.

The function $\Psi_k$ is convex. Each tangent to its cubic branch, together with its two linear branches, is a global supporting line. Rational lower bounds for cube roots and an upper bound for $\omega$ turn~\eqref{eq:HJB} into a sufficient affine inequality on a rectangle of $(X_1,X_2)$ values. Its minimum is at a corner. If no available supporting line proves the inequality on a rectangle, that rectangle is bisected. Every final comparison is between exact rational numbers. The complete calculation uses 8884 such rectangles and gives
\[
 \frac{3\sqrt3\pi}{\Lambda_0}w(1,1,s_{\rm top})>0.0879293>0.0872.
\]

The checker uses
\[
 3.14159265358979<\pi<3.14159265358980.
\]
These inequalities follow by substituting the alternating arctangent series into Machin's identity $\pi=16\arctan(1/5)-4\arctan(1/239)$ and bounding each remainder by its first omitted term. Integer square and cube roots similarly give rational enclosures for the algebraic constants. Floating-point estimates are used only to choose candidate supporting lines; the supporting-line condition and the final corner comparisons are then checked exactly.
Binary64 entries in the certificate are interpreted as exact dyadic rationals. The stored width and energy are slightly above $0.9$ and $37.7$. This is conservative: excess over a larger width is no larger, and excluding two spectral bottoms below a larger energy implies the stated bound at $37.7$. The finite test uses the exact stored parameters, rather than identifying them with the displayed decimals.

\subsection{The small-window certificates}
The barrier is $\rho=K\cot\Theta$, with $K=20$ and \emph{$\Theta$} continuous and affine on each triangle. The function $\rho$ itself is not piecewise affine. The angular conditions are
\[
 \Theta_s\ge\frac{\Lambda_0\sin^2\Theta+K^2\cos^2\Theta/\widetilde P^2}{K},
 \qquad
 \Theta_\ell\ge-\frac{2\sigma}{K}\sin^2\Theta,
 \qquad 0<\Theta<\pi.
\]
They imply~\eqref{eq:rho}. On each triangle, interval bounds for sine and cosine and a lower bound for $\widetilde P$ verify these inequalities throughout the triangle.

For the arc bound, a circular segment has the parametrization
\[
 A(\phi)=\tau^2\frac{\phi-\sin\phi\cos\phi}{\sin^2\phi},
 \qquad L(\phi)=\frac{2\tau\phi}{\sin\phi}.
\]
On $0\le\phi\le\pi/2$, $\dd L/\dd A=\sin\phi/\tau$ increases, so tangents are lower bounds. The checker encloses their coefficients with Arb and takes cellwise lower bounds. The required segment areas satisfy $A\le m<\pi\tau^2/2$, so they stay in this convex range. Outside the physical constraint $s\ge2\tau\ell$, the segment area is replaced by its positive part; on physical paths the formula is unchanged.

The origin is treated separately. Its nodal value is exactly $\pi/2$, not a floating-point approximation. On the bottom edge of the first triangle, $|\cos\Theta|\le|\Theta-\pi/2|=\Theta_s s$, while $\widetilde P^2=4\pi s$. This removes the apparent singularity. Off that edge, the perimeter lower bound is at least $4\tau$ in the first cell. In other bottom-row cells, the lower bound is the smaller of the zero-window and positive-window bounds.

The Bessel constant uses the interval
\[
 2.40482555769577276<j_{0,1}<2.40482555769577277.
\]
Arb verifies opposite signs of $J_0$ at the endpoints. To identify the first zero, the alternating series gives $J_1(x)\ge x(1-x^2/8)/2>0$ up to the right endpoint; hence $J_0'=-J_1<0$. The checker retains the full interval, including decimal-conversion error.

Nine mass intervals cover $[0.03,0.2282]$, with endpoints
\[
 .03,\ .10,\ .15,\ .19,\ .21,\ .22,\ .225,\ .227,\ .228,\ .2282.
\]
On $[m_1,m_2]$, take $A_0=2m_2$, $w=\min\{0.241,m_2/\tau\}$ and a lower bound for $\sigma_{\Lambda_0}(m_1)$. Monotonicity of~\eqref{eq:sigma} makes this conservative for every $m'$ in the interval. Each rectangular grid has $2000\times48$ cells, split along their diagonals into 192000 triangles. Interval evaluation verifies both angular inequalities on every triangle. The calculation uses Arb for the trigonometric bounds, with no assumed error estimate for platform trigonometric functions.

\subsection{The large-window and neck certificates}
For the large-window test, the lower bounds on required mass are as follows. The displayed bounds have been rounded down.
\begin{center}
\begin{tabular}{ccc}\toprule
Mass interval & Required mass at least & Width bound\\\midrule
$[.108,.15]$ & $.210000$ & $.240$\\
$[.15,.18]$ & $.228444$ & $.240$\\
$[.18,.20]$ & $.228731$ & $.240$\\
$[.20,.21]$ & $.228954$ & $.240$\\
$[.21,.22]$ & $.229083$ & $.240$\\
$[.22,.224]$ & $.229230$ & $.240$\\
$[.224,.2265]$ & $.229296$ & $.240$\\
$[.2265,.2282]$ & $.229341$ & $.240$\\\bottomrule
\end{tabular}
\end{center}
Every bound exceeds the upper endpoint of its mass interval. Candidate slopes are chosen numerically, then independently tested by interval evaluation of~\eqref{eq:potential}. The cell containing the true $\pi/2$ includes the value $\sin^2\theta=1$. The grid node near the initial angle is a dyadic approximation $h$ to $\pi/2$. In every integral starting at $\pi/2$, the first cell contributes $p(\theta_{i+1}-\pi/2)$, evaluated with Arb, rather than $p(\theta_{i+1}-h)$. Feasibility of the slope is checked on the whole cell. Thus the initial angle in both $A_k$ and $B_k$ is the mathematical $\pi/2$, including its rounding enclosure. Interval-to-floating conversions used in bounds are rounded outwards and checked; the final positivity test uses Arb intervals.

For the neck estimate, the Riccati function starts with the disk solution
\[
 \zeta(s)=2\pi X\frac{J_1(X)}{J_0(X)},\qquad X=\sqrt{Ls/\pi},
\]
on $0\le s\le2m_0$. The checker verifies that there is no pole there. On each subsequent cell, a lower bound $I\ge d>0$ gives the explicit majorant
\[
 d\sqrt L\tan\left(\frac{\sqrt L}{d}(s-s_i)
                +\arctan\frac{\zeta(s_i)}{d\sqrt L}\right).
\]
Its upper interval endpoint initializes the next cell. Upward jumps are allowed by the comparison principle. There are 100000 cells.

All comparisons use outward-rounded intervals or exact rationals. An interval containing zero is inconclusive and is never counted as a positive lower bound. As implementation checks, the calculation is repeated with deliberately unsupported stronger bounds and with one corrupted barrier node; these tests are rejected. Such rejection tests the calculation, not the truth of the stronger mathematical assertion.

\appendix
\section{Covering the large-window trajectories}\label{app:stairs}
This appendix explains the global step after the cellwise potential inequality~\eqref{eq:potential}. Fix a floor $\theta_b=0.25$ and levels
\[
 \theta_R=t_0<t_1<\cdots<t_K,
\]
with the last level close to the upper stationary angle for the smallest allowed width $w$. A trajectory satisfying the terminal condition falls into one of three classes.

If it falls below $\theta_b$, it must later return to $\theta_R$. Use a potential with $\nu=w$ on that return segment and the elementary bound $b\ge w$ on the remainder. This gives
\[
 D=w\tau+\int_{\theta_b}^{\theta_R}p_{\max}(t;w)\dd t.
\]
If it reaches $t_K$, apply the same potential before its first arrival there to obtain
\[
 A_K=w\tau+\int_{\pi/2}^{t_K}p_{\max}(t;w)\dd t.
\]
These estimates need no lower bound on the earlier minimum angle: with $\nu=w$, the slope zero is admissible outside the stated interval.

Otherwise the maximum angle lies in some $[t_k,t_{k+1})$, and, after the first class has been excluded, the whole path stays in $[\theta_b,t_{k+1}]$. First arrival at $t_k$ gives
\[
 A_k=w\tau+\int_{\pi/2}^{t_k}p_{\max}(t;w)\dd t.
\]
A second estimate uses a common time coefficient $\nu_k$ and two potentials, one before and one after that first arrival:
\[
 B_k=\nu_k\tau+
       \int_{\pi/2}^{t_k}p_{\max}(t;\nu_k)\dd t
       -\int_{\theta_R}^{t_k}p_{\min}(t;\nu_k)\dd t.
\]
Both potentials satisfy~\eqref{eq:potential} on every visited cell. For the second one, the selected lower slopes are nonnegative on $[\theta_R,t_{k+1}]$; replacing the actual terminal angle, which is at least $\theta_R$, by $\theta_R$ can only weaken the bound. The two first-arrival times cancel because their time coefficient is the same. A value $\nu_k<w$ merely weakens this estimate and does not invalidate it.

The resulting universal lower bound is
\[
 \min\left\{D,A_K,\min_{0\le k<K}\max(A_k,B_k)\right\}.
\]
The checker verifies every slope over its entire angular cell and evaluates every term with interval arithmetic. It uses 60 levels. Thus turning trajectories and waiting near a stationary angle are included; monotonicity of the trajectory is never assumed.

\clearpage
\begin{thebibliography}{99}
\bibitem{Henrot} A. Henrot, \emph{Extremum Problems for Eigenvalues of Elliptic Operators}, Frontiers in Mathematics, Birkh\"auser, Basel, 2006. \href{https://doi.org/10.1007/3-7643-7706-2}{doi:10.1007/3-7643-7706-2}.
\bibitem{Polya} G. P\'olya, On the eigenvalues of vibrating membranes, \emph{Proc. London Math. Soc.} (3) \textbf{11} (1961), 419--433. \href{https://doi.org/10.1112/plms/s3-11.1.419}{doi:10.1112/plms/s3-11.1.419}.
\bibitem{Balls} N. Filonov, M. Levitin, I. Polterovich and D. A. Sher, P\'olya's conjecture for Euclidean balls, \emph{Invent. Math.} \textbf{234} (2023), 129--169. \href{https://doi.org/10.1007/s00222-023-01198-1}{doi:10.1007/s00222-023-01198-1}.
\bibitem{Annuli} N. Filonov, M. Levitin, I. Polterovich and D. A. Sher, P\'olya's conjecture for Dirichlet eigenvalues of annuli, \emph{J. London Math. Soc.} \textbf{113} (2026), e70425. \href{https://doi.org/10.1112/jlms.70425}{doi:10.1112/jlms.70425}.
\bibitem{Speciel} R. Speciel, Planar Dirichlet eigenvalue bounds from a trace identity, preprint (2026), \href{https://arxiv.org/abs/2610.05636v1}{arXiv:2610.05636v1}.
\bibitem{TangWang} Q. Tang and Z. Wang, Strict Dirichlet P\'olya inequality for simply connected planar domains, preprint (2026), \href{https://arxiv.org/abs/2610.13139v1}{arXiv:2610.13139v1}.
\bibitem{Talenti} G. Talenti, Elliptic equations and rearrangements, \emph{Ann. Scuola Norm. Sup. Pisa Cl. Sci.} (4) \textbf{3} (1976), 697--718. \url{https://numdam.org/item/ASNSP_1976_4_3_4_697_0/}.
\bibitem{BrothersZiemer} J. E. Brothers and W. P. Ziemer, Minimal rearrangements of Sobolev functions, \emph{J. Reine Angew. Math.} \textbf{384} (1988), 153--179. \href{https://doi.org/10.1515/crll.1988.384.153}{doi:10.1515/crll.1988.384.153}.
\bibitem{BrockSolynin} F. Brock and A. Yu. Solynin, An approach to symmetrization via polarization, \emph{Trans. Amer. Math. Soc.} \textbf{352} (2000), 1759--1796. \href{https://doi.org/10.1090/S0002-9947-99-02558-1}{doi:10.1090/S0002-9947-99-02558-1}.
\bibitem{Zettl} A. Zettl, \emph{Sturm--Liouville Theory}, Mathematical Surveys and Monographs, vol. 121, American Mathematical Society, Providence, RI, 2005. \href{https://doi.org/10.1090/surv/121}{doi:10.1090/surv/121}.
\end{thebibliography}

\end{document}
